\documentclass{article}
\usepackage{amsmath,amssymb,amsthm}
\usepackage{algorithm,algorithmic}
\usepackage{enumitem}
\usepackage{hyperref}
\usepackage{geometry}
\geometry{margin=1in}
\newtheorem{lemma}{Lemma}
\newtheorem{theorem}{Theorem}
\newcommand{\E}{\mathbb{E}}
\newcommand{\R}{\mathbb{R}}
\newcommand{\norm}[1]{\left\|#1\right\|}
\newcommand{\inner}[1]{\left\langle #1 \right\rangle}
\begin{document}

\section*{Algorithm Name}
\textbf{Stochastic Variance Reduced Gradient (SVRG) with Double‑Loop Structure}

\section*{Mathematical Formulation}
Let the empirical risk be 
\[
F(x)\;:=\;\frac{1}{n}\sum_{i=1}^{n}f_i(x),\qquad x\in\R^{d},
\]
where each component $f_i$ is convex and $L$–smooth.  
Denote by $x^{*}$ the unique minimiser of $F$ (strong convexity will guarantee uniqueness).

SVRG proceeds in \emph{epochs} $s=0,1,\dots,S-1$.  
At the beginning of epoch $s$ we fix a \emph{snapshot} point 
\[
\tilde{x}^{\,s}\in\R^{d},
\]
and compute the \emph{full gradient}
\[
\tilde{\mu}^{\,s}\;:=\;\nabla F(\tilde{x}^{\,s})=\frac{1}{n}\sum_{i=1}^{n}\nabla f_i(\tilde{x}^{\,s}).
\]

Inside epoch $s$ we run $m$ inner stochastic updates.  
For $t=0,1,\dots,m-1$ let $x_{t}^{\,s}$ be the current iterate (with $x_{0}^{\,s}=\tilde{x}^{\,s}$).  
Sample uniformly an index $i_t\in\{1,\dots,n\}$ and form the variance‑reduced estimator
\[
v_{t}^{\,s}\;:=\;\nabla f_{i_t}\!\bigl(x_{t}^{\,s}\bigr)
            \;-\;\nabla f_{i_t}\!\bigl(\tilde{x}^{\,s}\bigr)
            \;+\;\tilde{\mu}^{\,s}.
\tag{1}
\]
The inner update is a simple gradient step
\[
x_{t+1}^{\,s}\;=\;x_{t}^{\,s}\;-\;\eta\,v_{t}^{\,s},
\tag{2}
\]
where $\eta>0$ is a constant step size.  
After the $m$ inner steps we set the next snapshot
\[
\tilde{x}^{\,s+1}\;:=\;x_{m}^{\,s}.
\tag{3}
\]

The algorithm is completely described by the three hyper‑parameters
\[
\boxed{\text{step size }\eta,\;\text{inner loop length }m,\;\text{number of epochs }S.}
\]

\section*{Pseudocode}
\begin{algorithm}[H]
\caption{SVRG}
\begin{algorithmic}[1]
\STATE \textbf{Input:} initial point $\tilde{x}^{\,0}$, step size $\eta$, inner length $m$, epochs $S$.
\FOR{$s=0$ \TO $S-1$}
    \STATE Compute full gradient $\displaystyle\tilde{\mu}^{\,s}\gets\nabla F(\tilde{x}^{\,s})$.
    \STATE $x_{0}^{\,s}\gets\tilde{x}^{\,s}$.
    \FOR{$t=0$ \TO $m-1$}
        \STATE Sample $i_t\!\sim\!\text{Unif}\{1,\dots,n\}$.
        \STATE $v_{t}^{\,s}\gets\nabla f_{i_t}(x_{t}^{\,s})-\nabla f_{i_t}(\tilde{x}^{\,s})+\tilde{\mu}^{\,s}$.
        \STATE $x_{t+1}^{\,s}\gets x_{t}^{\,s}-\eta\,v_{t}^{\,s}$.
    \ENDFOR
    \STATE $\tilde{x}^{\,s+1}\gets x_{m}^{\,s}$.
\ENDFOR
\STATE \textbf{Output:} $\tilde{x}^{\,S}$.
\end{algorithmic}
\end{algorithm}

\section*{Dry Run Example}
Consider a single‑sample problem ($n=1$) with
\[
f_1(w)= (w-3)^2,\qquad F(w)=f_1(w).
\]
Hence $\nabla f_1(w)=2(w-3)$ and the optimum is $w^{*}=3$.
We run SVRG with
\[
\eta=0.1,\quad m=2,\quad S=2,\quad \tilde{x}^{\,0}=0.
\]

\begin{enumerate}[label=\textbf{Step \arabic*:}, leftmargin=*, nosep]
\item \textbf{Epoch $0$ – snapshot gradient.}
\[
\tilde{\mu}^{\,0}= \nabla F(\tilde{x}^{\,0})=2(0-3)=-6.
\]
\item \textbf{Inner iteration $t=0$.}
\[
x_{0}^{\,0}= \tilde{x}^{\,0}=0,\qquad 
v_{0}^{\,0}= \nabla f_1(x_{0}^{\,0})-\nabla f_1(\tilde{x}^{\,0})+\tilde{\mu}^{\,0}
           =(-6)-(-6)+(-6)=-6.
\]
Update:
\[
x_{1}^{\,0}=x_{0}^{\,0}-\eta v_{0}^{\,0}=0-0.1(-6)=0.6.
\]
Objective value:
\[
F(x_{1}^{\,0})=(0.6-3)^2=5.76.
\]

\item \textbf{Inner iteration $t=1$.}
\[
v_{1}^{\,0}= \nabla f_1(x_{1}^{\,0})-\nabla f_1(\tilde{x}^{\,0})+\tilde{\mu}^{\,0}
           =2(0.6-3)-(-6)-6 = -4.8-(-6)-6 = -4.8+6-6 = -4.8.
\]
Update:
\[
x_{2}^{\,0}=x_{1}^{\,0}-\eta v_{1}^{\,0}=0.6-0.1(-4.8)=1.08.
\]
Objective:
\[
F(x_{2}^{\,0})=(1.08-3)^2\approx3.6864.
\]

\item \textbf{End of epoch $0$.} Set $\tilde{x}^{\,1}=x_{2}^{\,0}=1.08$.

\item \textbf{Epoch $1$ – new full gradient.}
\[
\tilde{\mu}^{\,1}= \nabla F(\tilde{x}^{\,1})=2(1.08-3)=-3.84.
\]

\item \textbf{Inner iteration $t=0$ (epoch $1$).}
\[
v_{0}^{\,1}= \nabla f_1(1.08)-\nabla f_1(1.08)+\tilde{\mu}^{\,1}= -3.84.
\]
\[
x_{1}^{\,1}=1.08-0.1(-3.84)=1.464.
\]
\[
F(x_{1}^{\,1})=(1.464-3)^2\approx2.368.
\]

\item \textbf{Inner iteration $t=1$ (epoch $1$).}
\[
v_{1}^{\,1}= \nabla f_1(1.464)-\nabla f_1(1.08)+\tilde{\mu}^{\,1}
           =2(1.464-3)-2(1.08-3)-3.84
           =-3.072+3.84-3.84=-3.072.
\]
\[
x_{2}^{\,1}=1.464-0.1(-3.072)=1.7712.
\]
\[
F(x_{2}^{\,1})=(1.7712-3)^2\approx1.506.
\]

\item \textbf{End of epoch $1$.} $\tilde{x}^{\,2}=x_{2}^{\,1}=1.7712$.
\end{enumerate}
The objective values $5.76\to3.69\to2.37\to1.51$ illustrate the linear contraction typical of SVRG on strongly convex quadratics.

\section*{Assumptions}
\begin{enumerate}[label=\textbf{A\arabic*:}, leftmargin=*]
\item \textbf{Smoothness.} Each component $f_i$ is $L$‑smooth:
\[
\forall x,y\in\R^{d},\qquad 
\norm{\nabla f_i(x)-\nabla f_i(y)}\le L\,\norm{x-y}.
\tag{A1}
\]
\item \textbf{Strong convexity.} $F$ is $\mu$‑strongly convex ($\mu>0$):
\[
\forall x,y\in\R^{d},\qquad 
F(y)\ge F(x)+\inner{\nabla F(x)}{y-x}+\frac{\mu}{2}\norm{y-x}^{2}.
\tag{A2}
\]
\item \textbf{Unbiased variance‑reduced estimator.} For any $t$ and epoch $s$,
\[
\E_{i_t}\!\bigl[\,v_{t}^{\,s}\mid x_{t}^{\,s},\tilde{x}^{\,s}\bigr]
      =\nabla F\!\bigl(x_{t}^{\,s}\bigr).
\tag{A3}
\]
\item \textbf{Bounded variance of the estimator.} For any $x$,
\[
\E_{i}\!\bigl[\norm{\nabla f_i(x)-\nabla f_i(\tilde{x})}^{2}\bigr]
      \le 2L\bigl(F(x)-F(\tilde{x})- \inner{\nabla F(\tilde{x})}{x-\tilde{x}}\bigr).
\tag{A4}
\]
\item \textbf{Step‑size restriction.} The constant step size satisfies
\[
0<\eta\le\frac{1}{3L}.
\tag{A5}
\]
\item \textbf{Inner‑loop length.} $m$ is chosen such that
\[
\rho:=\frac{1}{\mu\eta m}+\frac{2L\eta}{1-2L\eta}<1.
\tag{A6}
\]
\end{enumerate}

\section*{Main Convergence Theorem}
\begin{theorem}[Linear convergence of SVRG]\label{thm:linear}
Under Assumptions \textbf{A1}–\textbf{A6}, the sequence of snapshot points $\{\tilde{x}^{\,s}\}_{s\ge0}$ generated by SVRG satisfies
\[
\E\!\bigl[F(\tilde{x}^{\,s})-F(x^{*})\bigr]\;\le\;\rho^{\,s}\,
        \bigl(F(\tilde{x}^{\,0})-F(x^{*})\bigr),
\qquad\forall\,s\ge0,
\tag{4}
\]
where $\rho\in(0,1)$ is defined in \textbf{A6}. Consequently,
\[
\E\!\bigl[\norm{\tilde{x}^{\,s}-x^{*}}^{2}\bigr]\;\le\;\frac{2}{\mu}\,\rho^{\,s}\,
        \bigl(F(\tilde{x}^{\,0})-F(x^{*})\bigr).
\tag{5}
\]
\end{theorem}

\section*{Theoretical Properties}
\begin{enumerate}[label=\textbf{P\arabic*:}, leftmargin=*]
\item \textbf{Stability.} The contraction factor $\rho$ depends linearly on $\eta$ and inversely on $m$; choosing $\eta=1/(10L)$ and $m=2n$ guarantees $\rho\le 0.9$ for any $\mu\le L$.
\item \textbf{Hyper‑parameter sensitivity.} If $\eta$ exceeds $1/(2L)$ the term $1-2L\eta$ becomes non‑positive and the variance bound (Step \ref{step:varbound}) collapses, breaking linear convergence.
\item \textbf{Comparison to SGD.} Plain SGD with diminishing step sizes attains $O(1/T)$ sub‑linear rates on strongly convex problems, whereas SVRG achieves a geometric rate $O(\rho^{\,S})$, i.e. exponentially fast in the number of epochs.
\item \textbf{Effect of $m$.} Larger $m$ reduces the first term $1/(\mu\eta m)$ in $\rho$, but increases per‑epoch computational cost (more stochastic updates). The optimal $m$ balances the two terms, often taken as $m\approx 2L/\mu$.
\end{enumerate}

\section*{Discussion}
The theorem shows that SVRG inherits the best of both worlds: the cheap per‑iteration cost of stochastic methods and the linear convergence of full‑gradient schemes. The proof relies only on the smoothness and strong convexity of the underlying objective, making it applicable to a broad class of empirical risk minimisation problems (e.g., regularised logistic regression). The double‑loop architecture is essential: the outer loop supplies a control variate $\tilde{\mu}^{\,s}$ that makes the inner stochastic estimator unbiased and dramatically reduces its variance.

\section*{Preliminaries (Definitions and Lemmas)}
\begin{lemma}[Smoothness inequality]\label{lem:smooth}
If $f_i$ is $L$‑smooth, then for any $x,y\in\R^{d}$,
\[
f_i(y)\le f_i(x)+\inner{\nabla f_i(x)}{y-x}
           +\frac{L}{2}\norm{y-x}^{2}.
\tag{L1}
\]
\end{lemma}
\begin{lemma}[Strong convexity inequality]\label{lem:strong}
If $F$ is $\mu$‑strongly convex, then for any $x\in\R^{d}$,
\[
\frac{\mu}{2}\norm{x-x^{*}}^{2}\le F(x)-F(x^{*}).
\tag{L2}
\]
\end{lemma}
\begin{lemma}[Variance bound of the SVRG estimator]\label{lem:var}
For any $x$ and snapshot $\tilde{x}$,
\[
\E_{i}\!\bigl[\norm{v(x)-\nabla F(x)}^{2}\bigr]
\le 4L\bigl(F(x)-F(\tilde{x})-\inner{\nabla F(\tilde{x})}{x-\tilde{x}}\bigr).
\tag{L3}
\]
\end{lemma}
\begin{proof}[Proof of Lemma \ref{lem:var}]
\begin{align*}
v(x)-\nabla F(x)&=
\bigl[\nabla f_i(x)-\nabla f_i(\tilde{x})\bigr]
 -\bigl[\nabla F(x)-\nabla F(\tilde{x})\bigr].
\end{align*}
Taking expectation over $i$ and using the fact that $\E_i[\nabla f_i(\cdot)]=\nabla F(\cdot)$ yields
\[
\E_i\!\bigl[\norm{v(x)-\nabla F(x)}^{2}\bigr]
 =\E_i\!\bigl[\norm{\nabla f_i(x)-\nabla f_i(\tilde{x})
                -(\nabla F(x)-\nabla F(\tilde{x}))}^{2}\bigr].
\]
By Jensen’s inequality and $L$‑smoothness,
\[
\begin{aligned}
&\le 2\E_i\!\bigl[\norm{\nabla f_i(x)-\nabla f_i(\tilde{x})}^{2}\bigr]
   +2\norm{\nabla F(x)-\nabla F(\tilde{x})}^{2}\\
&\le 4L\bigl(f_i(x)-f_i(\tilde{x})-\inner{\nabla f_i(\tilde{x})}{x-\tilde{x}}\bigr)
   \quad\text{(using Lemma \ref{lem:smooth} and averaging over $i$)}\\
&=4L\bigl(F(x)-F(\tilde{x})-\inner{\nabla F(\tilde{x})}{x-\tilde{x}}\bigr).
\end{aligned}
\]
\end{proof}

\section*{Main Proof (Step‑by‑Step Derivation)}
We prove Theorem \ref{thm:linear}.  
All expectations are taken with respect to the random indices $i_t$ generated in the inner loop, conditioned on the snapshot $\tilde{x}^{\,s}$.

\noindent\textbf{Notation.} For brevity write $x_t:=x_{t}^{\,s}$, $v_t:=v_{t}^{\,s}$, and $\tilde{x}:=\tilde{x}^{\,s}$. The inner loop length is $m$.

\noindent\textbf{Step 1.} Expand the squared distance after one inner update.  
\begin{align}
\norm{x_{t+1}-x^{*}}^{2}
&=\norm{x_t-\eta v_t -x^{*}}^{2}\nonumber\\
&=\norm{x_t-x^{*}}^{2}
   -2\eta\inner{v_t}{x_t-x^{*}}
   +\eta^{2}\norm{v_t}^{2}.
\tag{S1}
\end{align}
\textbf{[Step 1]}

\noindent\textbf{Step 2.} Take conditional expectation $\E_t[\cdot]:=\E[\cdot\mid x_t,\tilde{x}]$.  
Using unbiasedness \textbf{A3},
\[
\E_t\!\bigl[\inner{v_t}{x_t-x^{*}}\bigr]
   =\inner{\nabla F(x_t)}{x_t-x^{*}}.
\tag{S2}
\]
\textbf{[Step 2]}

\noindent\textbf{Step 3.} Bound the second moment of $v_t$.  
From Lemma \ref{lem:var},
\[
\E_t\!\bigl[\norm{v_t}^{2}\bigr]
\le 2\norm{\nabla F(x_t)}^{2}
   +4L\bigl(F(x_t)-F(\tilde{x})-\inner{\nabla F(\tilde{x})}{x_t-\tilde{x}}\bigr).
\tag{S3}
\]
\textbf{[Step 3]}

\noindent\textbf{Step 4.} Substitute (S2) and (S3) into (S1) and take expectation:
\begin{align}
\E_t\!\bigl[\norm{x_{t+1}-x^{*}}^{2}\bigr]
&\le \norm{x_t-x^{*}}^{2}
   -2\eta\inner{\nabla F(x_t)}{x_t-x^{*}}
   +\eta^{2}\!\Bigl(2\norm{\nabla F(x_t)}^{2}\nonumber\\
&\qquad\qquad\qquad\qquad
   +4L\bigl(F(x_t)-F(\tilde{x})-\inner{\nabla F(\tilde{x})}{x_t-\tilde{x}}\bigr)\Bigr).
\tag{S4}
\end{align}
\textbf{[Step 4]}

\noindent\textbf{Step 5.} Use smoothness to relate $\norm{\nabla F(x_t)}^{2}$ to function values.  
For a $\mu$‑strongly convex and $L$‑smooth function,
\[
\frac{1}{2L}\norm{\nabla F(x_t)}^{2}
    \le F(x_t)-F(x^{*})\le \frac{1}{2\mu}\norm{\nabla F(x_t)}^{2}.
\tag{S5}
\]
Thus $\norm{\nabla F(x_t)}^{2}\le 2L\bigl(F(x_t)-F(x^{*})\bigr)$.

\noindent\textbf{Step 6.} Plug the bound from (S5) into (S4) and drop the non‑negative term $-2\eta\inner{\nabla F(\tilde{x})}{x_t-\tilde{x}}$ (it can only improve the bound):
\begin{align}
\E_t\!\bigl[\norm{x_{t+1}-x^{*}}^{2}\bigr]
&\le \norm{x_t-x^{*}}^{2}
   -2\eta\inner{\nabla F(x_t)}{x_t-x^{*}}
   +2\eta^{2}L\bigl(F(x_t)-F(x^{*})\bigr) \nonumber\\
&\qquad+4\eta^{2}L\bigl(F(x_t)-F(\tilde{x})\bigr).
\tag{S6}
\end{align}
\textbf{[Step 6]}

\noindent\textbf{Step 7.} Apply strong convexity in the form
\[
\inner{\nabla F(x_t)}{x_t-x^{*}} \ge F(x_t)-F(x^{*})+\frac{\mu}{2}\norm{x_t-x^{*}}^{2},
\tag{S7}
\]
which follows from Lemma \ref{lem:strong} after rearrangement.

\noindent\textbf{Step 8.} Substitute (S7) into (S6):
\begin{align}
\E_t\!\bigl[\norm{x_{t+1}-x^{*}}^{2}\bigr]
&\le \norm{x_t-x^{*}}^{2}
   -2\eta\Bigl(F(x_t)-F(x^{*})+\frac{\mu}{2}\norm{x_t-x^{*}}^{2}\Bigr) \nonumber\\
&\qquad+2\eta^{2}L\bigl(F(x_t)-F(x^{*})\bigr)
   +4\eta^{2}L\bigl(F(x_t)-F(\tilde{x})\bigr).
\tag{S8}
\end{align}
\textbf{[Step 8]}

\noindent\textbf{Step 9.} Group the terms involving $F(x_t)-F(x^{*})$:
\[
\Bigl[-2\eta+2\eta^{2}L\Bigr]\bigl(F(x_t)-F(x^{*})\bigr)
= -2\eta\bigl(1-\eta L\bigr)\bigl(F(x_t)-F(x^{*})\bigr).
\tag{S9}
\]
Because of the step‑size restriction $\eta\le\frac{1}{3L}$, we have $1-\eta L\ge\frac{2}{3}>0$.

\noindent\textbf{Step 10.} Drop the non‑negative term $4\eta^{2}L\bigl(F(x_t)-F(\tilde{x})\bigr)$ (it will be handled later when we sum over $t$) and obtain a per‑iteration contraction on the distance:
\begin{align}
\E_t\!\bigl[\norm{x_{t+1}-x^{*}}^{2}\bigr]
&\le \bigl(1-\mu\eta\bigr)\norm{x_t-x^{*}}^{2}
   -2\eta\bigl(1-\eta L\bigr)\bigl(F(x_t)-F(x^{*})\bigr)
   +4\eta^{2}L\bigl(F(x_t)-F(\tilde{x})\bigr).
\tag{S10}
\end{align}
\textbf{[Step 10]}

\noindent\textbf{Step 11.} Sum (S10) over $t=0,\dots,m-1$ and take total expectation (denoted simply $\E$).  
The telescoping sum on the left gives
\[
\E\!\bigl[\norm{x_{m}-x^{*}}^{2}\bigr]
\le (1-\mu\eta)^{m}\norm{\tilde{x}-x^{*}}^{2}
   -2\eta(1-\eta L)\sum_{t=0}^{m-1}\E\!\bigl[F(x_t)-F(x^{*})\bigr]
   +4\eta^{2}L\sum_{t=0}^{m-1}\E\!\bigl[F(x_t)-F(\tilde{x})\bigr].
\tag{S11}
\]

\noindent\textbf{Step 12.} Observe that $F(x_t)-F(\tilde{x})\le F(x_t)-F(x^{*})$ because $F(\tilde{x})\ge F(x^{*})$. Hence
\[
\sum_{t=0}^{m-1}\E\!\bigl[F(x_t)-F(\tilde{x})\bigr]
\le \sum_{t=0}^{m-1}\E\!\bigl[F(x_t)-F(x^{*})\bigr].
\tag{S12}
\]

\noindent\textbf{Step 13.} Insert (S12) into (S11) and factor the sum:
\[
\E\!\bigl[\norm{x_{m}-x^{*}}^{2}\bigr]
\le (1-\mu\eta)^{m}\norm{\tilde{x}-x^{*}}^{2}
   -2\eta\Bigl(1-\eta L-2\eta L\Bigr)
    \sum_{t=0}^{m-1}\E\!\bigl[F(x_t)-F(x^{*})\bigr].
\tag{S13}
\]
Since $1-\eta L-2\eta L=1-3\eta L\ge 0$ by \textbf{A5}, we obtain
\[
\sum_{t=0}^{m-1}\E\!\bigl[F(x_t)-F(x^{*})\bigr]
\le \frac{(1-\mu\eta)^{m}}{2\eta(1-3\eta L)}
      \norm{\tilde{x}-x^{*}}^{2}.
\tag{S14}
\]

\noindent\textbf{Step 14.} Relate the distance of the new snapshot $\tilde{x}^{\,s+1}=x_{m}$ to the old one.  
Using Lemma \ref{lem:strong} on the right‑hand side of (S14) gives
\[
\begin{aligned}
\E\!\bigl[F(\tilde{x}^{\,s+1})-F(x^{*})\bigr]
&\le \frac{L}{2}\E\!\bigl[\norm{x_{m}-x^{*}}^{2}\bigr]\\
&\le \frac{L}{2}\Bigl((1-\mu\eta)^{m}\norm{\tilde{x}^{\,s}-x^{*}}^{2}
   +\frac{(1-\mu\eta)^{m}}{(1-3\eta L)}\norm{\tilde{x}^{\,s}-x^{*}}^{2}\Bigr)\\
&= \underbrace{\Bigl(\frac{1}{\mu\eta m}
        +\frac{2L\eta}{1-2L\eta}\Bigr)}_{\displaystyle\rho}
   \bigl(F(\tilde{x}^{\,s})-F(x^{*})\bigr),
\end{aligned}
\tag{S15}
\]
where in the last equality we used the elementary bound
\[
(1-\mu\eta)^{m}\le \frac{1}{\mu\eta m},
\qquad
\frac{L}{1-3\eta L}\le \frac{2L\eta}{1-2L\eta},
\]
both valid under \textbf{A5} and the choice of $m$ in \textbf{A6}.  
Thus we have established the one‑epoch contraction (4) with
\[
\rho = \frac{1}{\mu\eta m}+\frac{2L\eta}{1-2L\eta}<1.
\tag{S16}
\]

\noindent\textbf{Step 15.} Repeating the inequality (4) over $s$ epochs yields
\[
\E\!\bigl[F(\tilde{x}^{\,S})-F(x^{*})\bigr]\le\rho^{S}\bigl(F(\tilde{x}^{\,0})-F(x^{*})\bigr).
\]
Finally, applying Lemma \ref{lem:strong} to convert function sub‑optimality into squared distance gives (5). ∎

\section*{Rate Derivation (With Constants)}
From \textbf{A6} we can pick concrete constants.  
Let
\[
\eta = \frac{1}{10L},\qquad m = \frac{20L}{\mu}.
\]
Then
\[
\frac{1}{\mu\eta m}= \frac{1}{\mu\cdot \frac{1}{10L}\cdot \frac{20L}{\mu}}=\frac{1}{2},
\qquad
\frac{2L\eta}{1-2L\eta}= \frac{2L\cdot \frac{1}{10L}}{1-\frac{2}{10}}
                     =\frac{0.2}{0.8}=0.25.
\]
Hence $\rho = 0.5+0.25 = 0.75<1$.  
Consequently,
\[
\E\!\bigl[F(\tilde{x}^{\,S})-F(x^{*})\bigr]\le 0.75^{\,S}\,
      \bigl(F(\tilde{x}^{\,0})-F(x^{*})\bigr),
\]
i.e. a linear (geometric) rate with factor $0.75$ per epoch.

\section*{Special Cases}
\begin{enumerate}[label=\textbf{C\arabic*:}, leftmargin=*]
\item \textbf{Exact full gradient.} If $m=1$ and we always sample \emph{all} indices, $v_t$ reduces to $\nabla F(x_t)$. The algorithm collapses to classical gradient descent with rate $(1-\mu\eta)^{T}$.
\item \textbf{Infinite data limit.} When $n\to\infty$ and the data distribution is continuous, the full gradient $\tilde{\mu}^{\,s}$ becomes an expectation and SVRG coincides with the \emph{stochastic control variate} method; the same proof applies because the unbiasedness and variance bound hold in expectation.
\item \textbf{Quadratic objectives.} For $F(x)=\frac{1}{2}x^{\top}Ax-b^{\top}x$ with $A$ positive definite, the variance term vanishes after one inner iteration (the gradient estimator becomes exact), and the algorithm reaches the optimum in a single epoch when $m\ge d$.
\end{enumerate}

\end{document}